\cleardoublepage
\chapter{ANALISIS MASALAH}
\label{chap:analisis-masalah}

Bab ini menganalisis kondisi pemantauan \textit{Water Treatment Plant} (WTP) saat ini, 
ruang lingkup \textit{frontend}, identifikasi masalah dan kebutuhan operator, 
serta dasar pemilihan solusi pengembangan \textit{dashboard} berbasis web.

\section{Analisis Kondisi Saat Ini}
Pada operasional WTP konvensional, pemantauan kualitas air masih mengandalkan 
instrumen fisik yang terpisah, pengamatan visual, dan pencatatan manual oleh operator. 
Pola ini mengharuskan operator untuk melakukan perbandingan data secara manual 
di lapangan. Hal tersebut tidak hanya meningkatkan beban kognitif, tetapi juga 
memperlambat respons operator terhadap perubahan parameter kritis seperti pH, kekeruhan, 
dan \textit{Total Dissolved Solids} (TDS). 

Keterlambatan dalam mengenali perubahan parameter ini berdampak langsung pada 
pengambilan keputusan terkait penyesuaian dosis bahan kimia yang akan diberikan, seperti 
kaporit dan \textit{Polyaluminium Chloride} (PAC). Selain itu, pencatatan riwayat 
dan log \textit{alarm} yang tidak terpusat dalam satu alur visual menyulitkan operato dalam proses 
penelusuran akar masalah (\textit{troubleshooting}) ketika terjadi anomali pada sistem. 

Saat ini, sistem otomasi yang lebih luas telah menyediakan lapisan sensor, mikrokontroler, 
dan \textit{backend}. Oleh karena itu, masalah yang menjadi fokus tugas akhir ini adalah 
belum tersedianya \textit{frontend} terpadu yang mampu menerjemahkan data mentah dari 
\textit{backend} menjadi informasi visual dan antarmuka interaktif yang sesuai dengan 
kebutuhan operasional operator.

\section{Batas Sistem dan Pemangku Kepentingan}

Aktor utama yang berinteraksi langsung dengan \textit{frontend} adalah operator WTP
yang bertanggung jawab melakukan pemantauan kondisi instalasi serta mendukung
proses operasional sistem otomasi dosis. Batas sistem \textit{frontend} dimulai dari
konsumsi layanan REST API yang disediakan oleh \textit{backend}, kemudian data yang
diterima diolah dan disajikan dalam bentuk antarmuka pengguna (\textit{User Interface})
yang diakses melalui \textit{browser} pada perangkat operator. Dengan demikian,
\textit{frontend} berperan sebagai media interaksi antara operator dengan sistem
otomasi tanpa berhubungan secara langsung dengan perangkat lapangan.

Ruang lingkup pengembangan \textit{frontend} tidak mencakup proses akuisisi data,
pemrosesan telemetri, maupun mekanisme pengendalian perangkat keras. Perangkat
keras sensor, modul ESP32, protokol komunikasi MQTT, arsitektur basis data,
pemrosesan telemetri, serta algoritma kendali dosis fisik berada di luar batas
pengembangan tugas akhir ini. Seluruh komponen tersebut diasumsikan telah
berfungsi dengan baik sehingga pengembangan difokuskan pada implementasi
\textit{frontend} yang mampu menyajikan informasi secara akurat, responsif, dan
mudah digunakan oleh operator.

\textit{Frontend} bertugas untuk mengelola dan menyajikan tiga kelompok data utama:
\begin{enumerate}
    \item Data telemetri, yang mencakup nilai pH, kekeruhan, dan TDS, baik secara
    terkini (\textit{real-time}) maupun historis;
    \item Data operasional, yang meliputi persentase level tangki bahan kimia,
    status pompa dosis, log alarm, dan nilai parameter \textit{offset} dosis; serta
    \item Data autentikasi, yang mengatur status sesi (\textit{session}) dan tingkat
    otorisasi profil operator.
\end{enumerate}

Berdasarkan ketiga kelompok data tersebut, \textit{frontend} menghasilkan berbagai
bentuk keluaran yang mendukung aktivitas pemantauan operator. Keluaran tersebut
meliputi ringkasan kondisi sistem dalam bentuk visual, grafik historis, indikator
status pompa dan tangki bahan kimia, notifikasi atau peringatan visual ketika
terjadi kondisi abnormal, serta fitur ekspor data pengukuran dalam format dokumen.
Selain menyajikan informasi, \textit{frontend} juga meneruskan permintaan interaksi
operator, seperti pengaturan \textit{offset} dosis maupun pengakuan (\textit{acknowledgement})
alarm, kepada \textit{backend} melalui REST API sehingga perubahan yang dilakukan
operator dapat diproses oleh sistem.

\section{Identifikasi Masalah Pengguna}
\label{sec:masalah-pengguna}
Berdasarkan analisis kondisi saat ini, masalah pengguna yang perlu diselesaikan
oleh antarmuka sistem dirumuskan sebagai berikut:
\begin{enumerate}
    \item Operator belum memiliki layar pemantauan terpusat untuk memahami 
    kondisi WTP secara holistik dan cepat.
    \item Operator kesulitan membedakan validitas data, terutama untuk mengenali 
    mana data yang akurat, sedang memuat (\textit{loading}), gagal diperoleh 
    (\textit{error}), atau kedaluwarsa akibat hilangnya koneksi sensor.
    \item Peninjauan tren kualitas air memerlukan visualisasi grafik dengan fitur 
    penyaringan waktu yang fleksibel dan kemampuan ekspor data untuk pelaporan.
    \item Interaksi kritis seperti penanganan alarm dan pengubahan \textit{offset} dosis 
    memerlukan proteksi hak akses dan umpan balik (\textit{feedback}) antarmuka yang jelas 
    agar tidak terjadi kesalahan operasional.
    \item Operator memiliki mobilitas tinggi sehingga \textit{dashboard} harus 
    dapat diakses dan menyesuaikan tampilan pada berbagai ukuran layar perangkat.
\end{enumerate}

\section{Kebutuhan Fungsional}
\label{sec:kebutuhan-fungsional}
Untuk menyelesaikan masalah pengguna yang telah diidentifikasi pada
Subbab~\ref{sec:masalah-pengguna}, kebutuhan fungsional \textit{frontend} dirumuskan pada
Tabel~\ref{tab:kebutuhan-fungsional}. Kolom "Masalah" menunjukkan
keterkaitan setiap kebutuhan dengan masalah pengguna yang diidentifikasi
pada Subbab~\ref{sec:masalah-pengguna}.

\begin{table}[H]
\centering
\caption{Kebutuhan fungsional \textit{frontend}}
\label{tab:kebutuhan-fungsional}
\begin{tabular}{|p{1.3cm}|p{10.5cm}|p{1.4cm}|}
\hline
\textbf{ID} & \textbf{Kebutuhan} & \textbf{Masalah} \\
\hline
KF-01 & Sistem menyediakan fitur \textit{login}, \textit{logout}, dan pembatasan akses halaman. & M2 \\
\hline
KF-02 & Sistem menampilkan ringkasan kepatuhan batas aman dan nilai parameter kualitas air terkini. & M1 \\
\hline
KF-03 & Sistem memvisualisasikan level tangki bahan kimia dan indikator status pompa dosis. & M4 \\
\hline
KF-04 & Sistem memberikan peringatan visual ketika data telemetri tidak tersedia atau telah melewati batas waktu pembaruan (kedaluwarsa). & M2, M4 \\
\hline
KF-05 & Sistem menyajikan grafik, statistik dasar, dan riwayat historis untuk parameter pH, kekeruhan, dan TDS. & M3 \\
\hline
KF-06 & Sistem menyediakan penyaringan rentang waktu data riwayat dan fitur ekspor ke dalam format CSV/XLSX. & M3 \\
\hline
KF-07 & Sistem menampilkan daftar alarm aktif dan historis terselesaikan, serta mendukung aksi pengakuan (\textit{acknowledge}) alarm. & M4 \\
\hline
KF-08 & Sistem menampilkan informasi pengaturan dosis saat ini dan menyediakan form untuk memberikan input \textit{offset} dosis yang valid. & M4 \\
\hline
KF-09 & Sistem menyediakan navigasi yang intuitif untuk perpindahan antartampilan \textit{dashboard}, detail sensor, riwayat, alarm, dan pengaturan pompa. & M1, M3 \\
\hline
\end{tabular}
\end{table}

Daftar kebutuhan fungsional akan menjadi acuan utama dan diselaraskan 
dengan struktur \textit{route}, desain komponen, serta kontrak API pada tahap implementasi.

\section{Kebutuhan Nonfungsional}
Kebutuhan nonfungsional mendefinisikan atribut kualitas dan batasan arsitektur 
dari \textit{frontend} yang akan dibangun. Kebutuhan tersebut terdiri atas:
\begin{enumerate}
    \item \textbf{Responsivitas UI:} Antarmuka harus bersifat responsif (\textit{responsive design}) 
    sehingga tetap proporsional dan mudah digunakan, baik pada layar \textit{desktop} di 
    ruang kontrol maupun pada perangkat bergerak saat operator berpatroli.
    \item \textbf{Keandalan Presentasi (UI/UX):} Antarmuka harus secara tangguh menangani 
    berbagai status pengambilan data, termasuk memberikan indikator \textit{loading}, 
    pesan \textit{error} yang jelas, serta penanda visual untuk sesi atau data sensor yang terputus.
    \item \textbf{Kemudahan Pemeliharaan (\textit{Maintainability}):} Basis kode harus dibangun 
    dengan arsitektur komponen yang dapat digunakan kembali (\textit{reusable components}) 
    dan menggunakan TypeScript untuk menjamin keamanan tipe data (\textit{type safety}) 
    sejak tahap kompilasi.
    \item \textbf{Kemudahan \textit{Deployment}:} Aplikasi harus dapat dibangun (\textit{build}) 
    dengan mudah untuk dipasang pada infrastruktur \textit{cloud}, serta mengadopsi 
    teknologi \textit{Progressive Web App} (PWA) agar dapat diinstal layaknya aplikasi lokal.
\end{enumerate}

\section{Analisis Pemilihan Solusi}

Persoalan yang teridentifikasi pada subbab sebelumnya, yaitu belum tersedianya antarmuka terpadu bagi operator, 
pada dasarnya dapat diselesaikan melalui beberapa pendekatan yang berbeda. Sebelum memutuskan untuk membangun 
antarmuka secara mandiri, dilakukan analisis terhadap solusi-solusi siap pakai (\textit{off-the-shelf}) yang 
telah tersedia, baik berupa perkakas \textit{business intelligence} maupun platform IoT. Analisis ini penting 
untuk memastikan bahwa keputusan membangun sistem sendiri didasarkan pada kebutuhan yang tidak dapat dipenuhi 
solusi eksisting, bukan sekadar preferensi teknis.

\subsection{Kriteria Evaluasi Solusi}
Kriteria evaluasi diturunkan langsung dari kebutuhan fungsional dan nonfungsional yang telah dirumuskan, yaitu:
\begin{enumerate}
    \item \textbf{Visualisasi telemetri dan riwayat (KF-02, KF-05, KF-06).} Kemampuan menyajikan nilai terkini, 
    grafik tren, riwayat data deret waktu, serta ekspor laporan.
    
    \item \textbf{Interaksi dua arah operator (KF-07, KF-08).} Kemampuan operator melakukan aksi yang mengubah 
    keadaan sistem, yaitu pengakuan alarm dan pengiriman nilai \textit{offset} dosis dengan validasi masukan.
    
    \item \textbf{Logika domain spesifik (KF-02, KF-04).} Kemampuan menyajikan evaluasi kepatuhan parameter 
    terhadap Permenkes No.~2 Tahun 2023 serta penanganan kondisi data kedaluwarsa sesuai risiko operasional.
    
    \item \textbf{Integrasi dengan \textit{backend} sistem.} Kemampuan mengonsumsi REST API yang dikembangkan 
    pada subsistem \textit{backend} tugas akhir ini, termasuk algoritma penentuan dosis yang menjadi inti sistem.
    
    \item \textbf{Kemandirian operasional.} Ketiadaan ketergantungan pada lisensi berbayar per pengguna maupun 
    layanan pihak ketiga yang dapat berubah kebijakannya.
\end{enumerate}

\subsection{Analisis Alternatif Solusi}
Berdasarkan kriteria tersebut, dianalisis empat kategori solusi siap pakai yang lazim digunakan pada sistem
pemantauan berbasis IoT, sebagaimana dirangkum pada Tabel~\ref{tab:komparasi-solusi}.

\begin{longtable}{p{2cm}p{2.5cm}p{4.2cm}p{5cm}}
\caption{Perbandingan alternatif solusi terhadap kebutuhan sistem} \label{tab:komparasi-solusi} \\
\toprule
\textbf{Solusi} & \textbf{Karakteristik} & \textbf{Kelebihan} & \textbf{Keterbatasan terhadap Kebutuhan} \\
\midrule
\endfirsthead

\multicolumn{4}{c}{\tablename\ \thetable\ Perbandingan alternatif solusi terhadap kebutuhan sistem (lanjutan)} \\
\toprule
\textbf{Solusi} & \textbf{Karakteristik} & \textbf{Kelebihan} & \textbf{Keterbatasan terhadap Kebutuhan} \\
\midrule
\endhead

\endfoot

\bottomrule
\endlastfoot

Power BI & Perkakas \textit{business intelligence} untuk pelaporan dan analisis data. &
Kemampuan visualisasi dan pelaporan yang matang; integrasi kuat dengan sumber data tabular; kurva belajar rendah bagi pengguna non-teknis. &
Berorientasi pada analisis data historis dan bersifat satu arah (hanya membaca). Tidak menyediakan mekanisme aksi operator seperti pengakuan alarm maupun pengiriman perintah dosis ke sistem. Memerlukan lisensi per pengguna. \\ \midrule
ThingSpeak & Platform IoT untuk pengumpulan dan visualisasi data sensor. &
Penyiapan cepat; integrasi langsung dengan perangkat IoT; tersedia paket gratis untuk penggunaan terbatas. &
Struktur kanal bersifat kaku dan kustomisasi antarmuka sangat terbatas. Paket gratis membatasi laju pengiriman pesan. Tidak memadai untuk alur kerja operator yang memerlukan aksi dan validasi masukan. \\ \midrule
Grafana & Platform visualisasi dan \textit{alerting} untuk data deret waktu. &
Sangat unggul dalam visualisasi telemetri; mekanisme \textit{alerting} berbasis ambang batas yang matang; bersifat sumber terbuka dan dapat dipasang mandiri. &
Berfokus pada lapisan visualisasi dan tidak menyimpan data sendiri, sehingga seluruh komponen lain harus dirakit terpisah. Aksi tulis-balik dari antarmuka menuju sistem tidak tersedia secara bawaan dan memerlukan pengaya tambahan. \\ \midrule
ThingsBoard & Platform IoT terpadu dengan manajemen perangkat, \textit{dashboard}, dan mesin aturan. &
Menyediakan \textit{dashboard}, manajemen alarm dengan pengakuan, serta widget kendali untuk mengirim perintah ke perangkat. Merupakan alternatif paling mendekati kebutuhan sistem. &
Membawa mesin aturan dan lapisan pemrosesan data sendiri, sehingga penggunaannya menuntut data sistem disesuaikan dengan model data platform. Hal ini menjadikan subsistem \textit{backend} beserta algoritma penentuan dosis yang dikembangkan pada tugas akhir ini menjadi redundan. Kustomisasi tampilan terbatas pada paradigma widget yang disediakan. \\
\end{longtable}


\subsection{Analisis Penentuan Solusi}
Berdasarkan Tabel~\ref{tab:komparasi-solusi}, dapat disimpulkan bahwa untuk kebutuhan pemantauan telemetri 
semata, solusi siap pakai seperti Grafana maupun ThingsBoard secara umum lebih efisien dibandingkan pengembangan 
antarmuka secara mandiri, karena tidak memerlukan penulisan kode dari awal. Namun, kebutuhan sistem pada tugas 
akhir ini melampaui pemantauan telemetri semata, sehingga terdapat tiga pertimbangan yang menjadi dasar keputusan 
untuk membangun antarmuka secara mandiri.

Pertama, sistem otomasi dosis WTP yang dikembangkan pada tugas akhir ini memiliki subsistem \textit{backend} 
tersendiri yang memuat logika penentuan dosis bahan kimia. Penggunaan platform IoT seperti ThingsBoard 
akan menuntut logika tersebut dipindahkan ke dalam mesin aturan platform, atau menuntut data dipaksakan mengikuti 
model data platform. Kedua kemungkinan tersebut menjadikan subsistem \textit{backend} yang menjadi inti sistem 
otomasi kehilangan perannya. Dengan mengembangkan antarmuka secara mandiri, integrasi dilakukan langsung terhadap 
REST API subsistem \textit{backend} sehingga pembagian tanggung jawab antarsubsistem tetap terjaga.

Kedua, sistem memerlukan penyajian logika domain yang tidak tersedia sebagai fungsi bawaan pada platform mana pun, 
yaitu evaluasi kepatuhan parameter kualitas air terhadap Permenkes No.~2 Tahun 2023 beserta representasi 
visualnya, serta penanganan kondisi data kedaluwarsa yang dirancang untuk mencegah operator mengambil keputusan 
berdasarkan pembacaan yang tidak mutakhir. Kedua kebutuhan tersebut menuntut kendali penuh atas perilaku dan 
tampilan antarmuka.

Ketiga, kebutuhan interaksi dua arah, yaitu pengakuan alarm dan pengiriman \textit{offset} dosis yang disertai 
validasi masukan, menuntut alur kerja yang dirancang khusus mengikuti prosedur operasional WTP. Meskipun sebagian 
platform menyediakan kemampuan serupa, alur kerjanya terikat pada paradigma widget yang disediakan platform dan 
tidak dapat disesuaikan secara bebas dengan hasil analisis kebutuhan operator.

Berdasarkan ketiga pertimbangan tersebut, tugas akhir ini mengimplementasikan \textit{frontend dashboard} yang 
dikembangkan secara mandiri sebagai aplikasi web responsif. Perlu ditegaskan bahwa keputusan ini tidak menyiratkan 
bahwa solusi siap pakai bersifat inferior. Untuk sistem yang kebutuhannya terbatas pada visualisasi telemetri dan 
peringatan berbasis ambang batas, solusi siap pakai justru merupakan pilihan yang lebih tepat. Keputusan 
pengembangan mandiri pada tugas akhir ini berlaku secara spesifik karena adanya kebutuhan integrasi dengan 
subsistem \textit{backend} internal, logika domain yang bersifat khusus, serta alur interaksi operator yang 
dirancang berdasarkan analisis kebutuhan.

\subsection{Analisis Pemilihan Tumpukan Teknologi (\textit{Tech Stack})}
Setelah platform aplikasi web ditetapkan, pemilihan tumpukan teknologi dilakukan secara spesifik untuk mengatasi kompleksitas 
data telemetri dan kontrol dosis.

\textbf{1. Next.js vs. Pendekatan React Tradisional} \\
Meskipun aplikasi dapat dibangun menggunakan pustaka React standar (seperti \textit{Create React App} atau Vite), 
\textit{framework} Next.js dipilih karena menyediakan arsitektur \textit{routing} bawaan dan manajemen tata letak 
yang lebih terstruktur. Pembagian arsitektur berbasis komponen pada Next.js mempermudah pengisolasian logika UI 
sehingga bagian grafik riwayat yang berat tidak akan memengaruhi performa komponen \textit{sidebar} navigasi. 

\textbf{2. TanStack Query vs. \textit{State Management} Global (Redux/Context)} \\
Data telemetri WTP (seperti pH dan kekeruhan) adalah \textit{server state} yang perubahannya terjadi di luar 
jangkauan aplikasi \textit{frontend}. Penggunaan \textit{state management} tradisional seperti Redux atau React 
Context murni akan membutuhkan penulisan kode tambahan yang sangat panjang (\textit{boilerplate}) hanya untuk 
menangani status \textit{loading} dan \textit{error}. TanStack Query dipilih karena didesain khusus untuk sinkronisasi 
data asinkron. Fitur \textit{background polling} yang dimilikinya memungkinkan \textit{dashboard} melakukan pembaruan 
data secara otomatis (misalnya setiap 5 detik) tanpa menyela interaksi operator, sebuah fitur kritis untuk sistem pemantauan 
\textit{real-time}.

\textbf{3. TypeScript vs. JavaScript Standar} \\
Kesalahan eksekusi antarmuka (\textit{runtime error}) pada sistem pengendalian bahan kimia (seperti \textit{offset} kaporit) 
memiliki risiko operasional yang tinggi. Oleh karena itu, TypeScript dipilih menggantikan JavaScript standar. TypeScript 
memberikan keamanan tipe data (\textit{type safety}) dengan mewajibkan pendefinisian kontrak antarmuka (\textit{interface contract}) 
untuk setiap struktur JSON yang dikirim maupun diterima dari REST API. Pendekatan ini memungkinkan sebagian besar \textit{bug} terkait 
manipulasi data terdeteksi sejak tahap kompilasi kode, bukan saat aplikasi sedang digunakan oleh operator.

\subsection{Penentuan Solusi Akhir}
Berdasarkan hasil analisis komparatif di atas, Tugas Akhir ini akan mengimplementasikan \textit{dashboard} pemantauan WTP 
menggunakan platform aplikasi web responsif berstandar PWA. Sistem ini akan dibangun di atas ekosistem Next.js dengan bahasa 
pemrograman TypeScript, serta memanfaatkan TanStack Query sebagai instrumen utama dalam melakukan orkestrasi \textit{server state} 
antara \textit{frontend} dan REST API \textit{backend}.